\documentclass[10pt,a4paper]{amsart}
\usepackage{amsmath,amssymb,amsthm,graphicx}
% Retained minimal preamble for the Stacks Project chapter sites-cohomology.tex (Cohomology on Sites).
% Verbatim excerpt of the chapter preamble: Xy-pic setup (native preamble lines 14-18) and the native
% theorem declarations with the mathematical macros (native preamble lines 154-217). Copied mechanically;
% nothing added, removed or redefined. The wrapper class and the amsmath/amssymb/amsthm/graphicx base
% packages are supplied by the assembler, so the class/externaldocument part of the native preamble is omitted.
% For commutative diagrams we use Xy-pic
\usepackage[all]{xy}

% We use 2cell for 2-commutative diagrams.
\xyoption{2cell}
% Theorem environments.
%
\theoremstyle{plain}
\newtheorem{theorem}[subsection]{Theorem}
\newtheorem{proposition}[subsection]{Proposition}
\newtheorem{lemma}[subsection]{Lemma}

\theoremstyle{definition}
\newtheorem{definition}[subsection]{Definition}
\newtheorem{example}[subsection]{Example}
\newtheorem{exercise}[subsection]{Exercise}
\newtheorem{situation}[subsection]{Situation}

\theoremstyle{remark}
\newtheorem{remark}[subsection]{Remark}
\newtheorem{remarks}[subsection]{Remarks}

\numberwithin{equation}{subsection}

% Macros
%
\def\lim{\mathop{\mathrm{lim}}\nolimits}
\def\colim{\mathop{\mathrm{colim}}\nolimits}
\def\Spec{\mathop{\mathrm{Spec}}}
\def\Hom{\mathop{\mathrm{Hom}}\nolimits}
\def\Ext{\mathop{\mathrm{Ext}}\nolimits}
\def\SheafHom{\mathop{\mathcal{H}\!\mathit{om}}\nolimits}
\def\SheafExt{\mathop{\mathcal{E}\!\mathit{xt}}\nolimits}
\def\Sch{\mathit{Sch}}
\def\Mor{\mathop{\mathrm{Mor}}\nolimits}
\def\Ob{\mathop{\mathrm{Ob}}\nolimits}
\def\Sh{\mathop{\mathit{Sh}}\nolimits}
\def\NL{\mathop{N\!L}\nolimits}
\def\CH{\mathop{\mathrm{CH}}\nolimits}
\def\proetale{{pro\text{-}\acute{e}tale}}
\def\etale{{\acute{e}tale}}
\def\QCoh{\mathit{QCoh}}
\def\Ker{\mathop{\mathrm{Ker}}}
\def\Im{\mathop{\mathrm{Im}}}
\def\Coker{\mathop{\mathrm{Coker}}}
\def\Coim{\mathop{\mathrm{Coim}}}

% Boxtimes
%
\DeclareMathSymbol{\boxtimes}{\mathbin}{AMSa}{"02}

%
% Macros for moduli stacks/spaces
%
\def\QCohstack{\mathcal{QC}\!\mathit{oh}}
\def\Cohstack{\mathcal{C}\!\mathit{oh}}
\def\Spacesstack{\mathcal{S}\!\mathit{paces}}
\def\Quotfunctor{\mathrm{Quot}}
\def\Hilbfunctor{\mathrm{Hilb}}
\def\Curvesstack{\mathcal{C}\!\mathit{urves}}
\def\Polarizedstack{\mathcal{P}\!\mathit{olarized}}
\def\Complexesstack{\mathcal{C}\!\mathit{omplexes}}
% \Pic is the operator that assigns to X its picard group, usage \Pic(X)
% \Picardstack_{X/B} denotes the Picard stack of X over B
% \Picardfunctor_{X/B} denotes the Picard functor of X over B
\def\Pic{\mathop{\mathrm{Pic}}\nolimits}
\def\Picardstack{\mathcal{P}\!\mathit{ic}}
\def\Picardfunctor{\mathrm{Pic}}
\def\Deformationcategory{\mathcal{D}\!\mathit{ef}}
\usepackage{cleveref}
\crefname{originalreference}{original reference}{original references}
\Crefname{originalreference}{Original reference}{Original references}
\makeatletter
\newlabel{lemma-derived-pullback-composition}{{0D5S}{source}{}{sourcepointer-1}{}}
\newlabel{lemma-derived-pullback-composition@cref}{{[originalreference][0][]0D5S}{[1][0][]source}{}{}{}}
\makeatother
\title{The derived pullback along a morphism of ringed topoi sends quasi-coherent objects to quasi-coherent objects}
\author{The Stacks Project Authors}
\date{stacks-project:sites-cohomology.tex:local-snapshot-20261011; GNU FDL 1.2 or later}
\begin{document}
\maketitle
\noindent\textit{Editorial note.} All delivered mathematical text is the author's own, copied from the pinned Stacks Project chapter file sites-cohomology.tex; no statement or proof step was rewritten, reordered or removed.\par\medskip
\noindent\textit{Editorial note.} The counted proof is the author's complete argument: the proof is written out in full: the commutative diagram of ringed topoi over a single object, the identification of both composites as R Gamma(U, K) tensor\textasciicircum{}L\_\{O(U)\} O'(U'), and the five-term chain of equalities verifying the cartesianity identity for an arbitrary morphism U' -> V'.\par\medskip
\noindent\textit{Editorial note.} The label lemma-cartesian-functorial names the counted statement, and the proof environment immediately following it in the source is delivered verbatim as its proof.\par\medskip
\noindent\textit{Editorial note.} The counted statement is carried by the author's own lemma environment in the chapter source.\par\medskip
\noindent\textit{Editorial note.} The statement and proof are copied from the pinned chapter file; the single cited input is kept as an external pointer with its canonical Stacks tag.\par\medskip
\noindent\textit{Editorial note.} The preamble is the chapter's own retained minimal preamble (sources/stacks-sites-cohomology/retained-preamble.tex), a verbatim excerpt of the native preamble: the Xy-pic setup and the native theorem declarations with the mathematical macros.\par\medskip
\noindent\textit{Allowed standard prerequisites (editorial).}\begin{itemize}
\item The morphism of ringed topoi g and the functor Lg\textasciicircum{}*. Same chapter, the standing morphism g : (Sh(C'), O') -> (Sh(C), O) of the section and the derived pullback Lg\textasciicircum{}*. Supplies the functor whose compatibility with quasi-coherence is asserted.
\item Quasi-coherent objects QC(O) and the cartesianity condition. Same chapter, the definition of QC(O) by the requirement that R Gamma(V, K) tensor\textasciicircum{}L\_\{O(V)\} O(U) -> R Gamma(U, K) be an isomorphism for every U -> V. Supplies the property verified for Lg\textasciicircum{}*K at the end of the proof.
\item Composition of derived pullbacks and the ringed-topos squares over single objects. Same chapter, the lemma on composition of derived pullbacks through the commutative diagram of ringed topoi built from an object U' of C'. Supplies the identity R Gamma(U', Lg\textasciicircum{}*K) = R Gamma(U, K) tensor\textasciicircum{}L\_\{O(U)\} O'(U').
\end{itemize}\medskip
\ifdefined\hypertarget\hypertarget{sourcepointer-1}{}\fi
\noindent\textit{Source reference.} Original reference 0D5S: sources/stacks-sites-cohomology/pinned-original-source.tex, source lines 3658--3658. Prerequisite pointer invoked at proof line(s) 11526; it names a standard fact of the same chapter that the counted argument uses, and every step of the counted proof built on it is retained verbatim in the delivered body.\par\medskip
\begin{lemma}
\label{lemma-cartesian-functorial}
Let $g : (\Sh(\mathcal{C}'), \mathcal{O}') \to (\Sh(\mathcal{C}), \mathcal{O})$
be as above. Then the functor $Lg^* : D(\mathcal{O}) \to D(\mathcal{O}')$
maps $\mathit{QC}(\mathcal{O})$ into $\mathit{QC}(\mathcal{O}')$.
\end{lemma}

\begin{proof}
Let $U' \in \Ob(\mathcal{C}')$ with image $U = u(U')$ in $\mathcal{C}$.
Let $pt$ denote the category with a single object and a single morphism.
Denote $(\Sh(pt), \mathcal{O}'(U'))$ and $(\Sh(pt), \mathcal{O}(U))$
the ringed topoi as indicated. Of course we identify the derived category
of modules on these ringed topoi with $D(\mathcal{O}'(U'))$ and
$D(\mathcal{O}(U))$. Then we have a commutative diagram of ringed topoi
$$
\xymatrix{
(\Sh(pt), \mathcal{O}'(U')) \ar[rr]_{U'} \ar[d] & &
(\Sh(\mathcal{C}'), \mathcal{O}') \ar[d]^g \\
(\Sh(pt), \mathcal{O}(U)) \ar[rr]^U & &
(\Sh(\mathcal{C}), \mathcal{O})
}
$$
Pullback along the lower horizontal morphism sends $K$ in $D(\mathcal{O})$
to $R\Gamma(U, K)$. Pullback by the left vertical arrow sends
$M$ to $M \otimes_{\mathcal{O}(U)}^\mathbf{L} \mathcal{O}'(U')$.
Going around the diagram either direction produces the same result
(Lemma \ref{lemma-derived-pullback-composition}) and hence we conclude
$$
R\Gamma(U', Lg^*K) =
R\Gamma(U, K) \otimes_{\mathcal{O}(U)}^\mathbf{L} \mathcal{O}'(U')
$$
Finally, let $f' : U' \to V'$ be a morphism in $\mathcal{C}'$ and
denote $f = u(f') : U = u(U') \to V = u(V')$ the image in $\mathcal{C}$.
If $K$ is in $\mathit{QC}(\mathcal{O})$ then we have
\begin{align*}
R\Gamma(V', Lg^*K) \otimes_{\mathcal{O}'(V')}^\mathbf{L} \mathcal{O}'(U')
& =
R\Gamma(V, K) \otimes_{\mathcal{O}(V)}^\mathbf{L} \mathcal{O}'(V')
\otimes_{\mathcal{O}'(V')}^\mathbf{L} \mathcal{O}'(U') \\
& =
R\Gamma(V, K) \otimes_{\mathcal{O}(V)}^\mathbf{L} \mathcal{O}'(U') \\
& =
R\Gamma(V, K) \otimes_{\mathcal{O}(V)}^\mathbf{L} \mathcal{O}(U)
\otimes_{\mathcal{O}(U)}^\mathbf{L} \mathcal{O}'(U') \\
& =
R\Gamma(U, K) \otimes_{\mathcal{O}(U)}^\mathbf{L} \mathcal{O}'(U') \\
& =
R\Gamma(U', Lg^*K)
\end{align*}
as desired. Here we have used the observation above both for $U'$ and $V'$.
\end{proof}

\end{document}
